% ============================================================================
% APPENDIX XVI: LIT-CRITIQUE - Critical Literature and Canonical Arguments
% ============================================================================
% Category: LIT (Literature/Critical Analysis)
% Language: English
% Version: 1.0
% Status: Complete
% Prerequisites: XV (LIT-HISTORY), X (FORMAL-FOUND), VII (GTC)
% ============================================================================

\begin{chapter}{Critical Literature: The Ten Canonical Arguments Against Complementarity}
\label{app:lit-critique}

% ============================================================================
% HEADER BLOCK
% ============================================================================
\begin{tcolorbox}[colback=gray!5!white,colframe=gray!75!black]
\small
\textbf{Category:} LIT (Literature/Critical Analysis) \\
\textbf{Code:} XVI \\
\textbf{Version:} 1.0 \\
\textbf{Status:} Complete \\
\textbf{Dependencies:} XV (LIT-HISTORY), X (FORMAL-FOUND), VII (GTC) \\
\textbf{Scientific Purpose:} Systematic documentation of objections and responses
\end{tcolorbox}

% ============================================================================
% CROSS-REFERENCE MAP
% ============================================================================
\begin{tcolorbox}[colback=blue!5!white,colframe=blue!75!black,title=\textbf{Cross-Reference Map}]
\small
\textbf{This appendix relates to:}
\begin{itemize}[noitemsep]
    \item \textbf{XV (LIT-HISTORY):} Why these objections explain historical non-development
    \item \textbf{X (FORMAL-FOUND):} Formal definitions that address conceptual objections
    \item \textbf{VII (GTC):} Complete theory that addresses all ten arguments
    \item \textbf{XI (FORMAL-MODELS):} Hierarchical models addressing scale objections
    \item \textbf{BBB (CORE-WHERE):} Estimation methods addressing identification
\end{itemize}
\end{tcolorbox}

% ============================================================================
% CHAPTER LINKAGE
% ============================================================================
\begin{tcolorbox}[colback=purple!5!white, colframe=purple!75!black,
    title=Chapter Linkage]

\textbf{Primary Chapter:} Chapter 4: Critical Assessment
\begin{itemize}[nosep]
\item Section 4.1: Canonical Objections $\leftrightarrow$ CRT.1-CRT.10
\item Section 4.2: Response Framework $\leftrightarrow$ CRT.11
\end{itemize}

\textbf{Secondary Chapters:}
\begin{itemize}[nosep]
\item Chapter 7: Mathematical Foundations --- formal resolutions
\item Chapter 8: Integration --- synthesis of responses
\end{itemize}

\end{tcolorbox}

% ============================================================================
% ABSTRACT
% ============================================================================
\begin{abstract}
This appendix systematically documents the critical literature on complementarity, distilling \textbf{ten canonical arguments} that recur across disciplines. These objections---concerning identification, model-dependence, scale, context, dominance, complexity, measurement noise, post-hoc narrativity, normativity, and conceptual vagueness---are not chaotic but \textit{structured}. This structure reveals that the critiques are \textbf{methodological, not ontological}: they explain why complementarity remained under-formalized, but do not refute its existence. The appendix presents 100 references (50 supporting, 50 contrasting) and demonstrates how the EBF $\gamma$-framework addresses each canonical objection. The central conclusion: \textit{the objections against complementarity are precisely what a formal theory must---and can---resolve.}
\end{abstract}

% -----------------------------------------------------------------------------
% APPENDIX SCOPE BOX
% -----------------------------------------------------------------------------

\begin{tcolorbox}[
    colback=orange!5!white,
    colframe=orange!75!black,
    title={\textbf{Appendix Scope: Ziel / In-Scope / Out-of-Scope / Constraints}},
    fonttitle=\bfseries
]

\textbf{Ziel (Objective):}
\begin{itemize}[nosep]
    \item Formalize and document the key concepts of Unknown
\end{itemize}

\textbf{In-Scope:}
\begin{itemize}[nosep]
    \item Core theoretical foundations
    \item Mathematical formalizations where applicable
    \item Integration with EBF framework
\end{itemize}

\textbf{Out-of-Scope (delegiert an andere Appendices):}
\begin{itemize}[nosep]
    \item Detailed empirical validation $\rightarrow$ METHOD appendices
    \item Domain-specific applications $\rightarrow$ DOMAIN appendices
    \item Complete literature review $\rightarrow$ LIT appendices
\end{itemize}

\textbf{Constraints (Anwendungsgrenzen):}
\begin{itemize}[nosep]
    \item Framework requires calibration to specific contexts
    \item Parameters may vary across domains
    \item Validity bounded by underlying assumptions
\end{itemize}

\textbf{Lieferobjekte (Deliverables):}
\begin{itemize}[nosep]
    \item Theoretical framework with formal definitions
    \item Integration points with other appendices
    \item Practical guidelines for application
\end{itemize}

\end{tcolorbox}

%%%%%%%%%%%%%%%%%%%%%%%%%%%%%%%%%%%%%%%%%%%%%%%%%%%%%%%%%%%%%%%%%%%%%%%%%%%%%%%
\section{The Fundamental Question}
\label{sec:crt-fundamental}
%%%%%%%%%%%%%%%%%%%%%%%%%%%%%%%%%%%%%%%%%%%%%%%%%%%%%%%%%%%%%%%%%%%%%%%%%%%%%%%

\begin{tcolorbox}[
    colback=red!5!white,
    colframe=red!75!black,
    title={\textbf{Fundamental Question}}
]
\textbf{How does unknown integrate with and contribute to the
Evidence-Based Framework for understanding behavioral outcomes?}

This appendix addresses the core challenge of formalizing behavioral principles
within a rigorous theoretical framework while maintaining practical applicability.
\end{tcolorbox}





% ============================================================================
% FUNDAMENTAL INSIGHT
% ============================================================================
\begin{tcolorbox}[colback=red!5!white,colframe=red!75!black,title=\textbf{Fundamental Insight}]
\begin{center}
\large
\textbf{The central objections to complementarity are not substantive but methodological---and therefore solvable.}
\end{center}

\vspace{0.5em}
The ten canonical arguments explain why complementarity was historically under-developed. They do not demonstrate that complementarity is false, irrelevant, or unmeasurable---only that it is \textit{difficult}. The $\gamma$-framework is the response.
\end{tcolorbox}

% ============================================================================
% QUICK REFERENCE
% ============================================================================
\begin{tcolorbox}[colback=yellow!5!white,colframe=yellow!75!black,title=\textbf{Quick Reference: The Ten Canonical Arguments}]
\small
\renewcommand{\arraystretch}{1.2}
\begin{tabular}{@{}clcc@{}}
\toprule
\textbf{\#} & \textbf{Argument} & \textbf{Strength} & \textbf{$\gamma$-Response} \\
\midrule
1 & Identification & Very High & Explicit decomposition \\
2 & Model-Dependence & High & Reference model explicit \\
3 & Scale & High & Hierarchical $\gamma^L$ \\
4 & Context & High & Conditional $\gamma|\Psi$ \\
5 & Dominance/Mass-Ratio & Medium-High & Weighted decomposition \\
6 & Complexity & Very High & Möbius structure \\
7 & Measurement Noise & Very High & Regularization \\
8 & Post-hoc Narrativity & Medium & Ex-ante design \\
9 & Normativity & Medium & $\gamma < 0$ explicit \\
10 & Conceptual Vagueness & Very High & Formal definition \\
\bottomrule
\end{tabular}
\end{tcolorbox}

% ============================================================================
% PART I: THE TEN CANONICAL ARGUMENTS
% ============================================================================
\section{The Ten Canonical Arguments Against Complementarity}
\label{sec:ten-arguments}

\subsection{Overview: Structure of Criticism}

The critical literature on complementarity, spanning economics, pharmacology, ecology, neuroscience, and machine learning, converges on remarkably similar objections. This convergence is scientifically significant: it reveals that the problems are \textit{structural}, not domain-specific.

\begin{figure}[h]
\centering
\begin{tikzpicture}[
    arg/.style={rectangle, draw, rounded corners, minimum width=2.8cm, minimum height=0.7cm, align=center, font=\small},
    cat/.style={ellipse, draw, fill=gray!20, minimum width=2cm, minimum height=0.6cm, align=center, font=\small}
]
% Categories
\node[cat] (epist) at (0,3) {Epistemic};
\node[cat] (method) at (4,3) {Methodological};
\node[cat] (pract) at (8,3) {Practical};

% Arguments
\node[arg, fill=blue!20] (a1) at (-1,1.5) {1. Identification};
\node[arg, fill=blue!20] (a2) at (1.5,1.5) {2. Model-Dep.};
\node[arg, fill=green!20] (a3) at (3,0.5) {3. Scale};
\node[arg, fill=green!20] (a4) at (5,1.5) {4. Context};
\node[arg, fill=green!20] (a5) at (5,0) {5. Dominance};
\node[arg, fill=orange!20] (a6) at (7,1.5) {6. Complexity};
\node[arg, fill=orange!20] (a7) at (9,1.5) {7. Noise};
\node[arg, fill=blue!20] (a8) at (0,0) {8. Post-hoc};
\node[arg, fill=blue!20] (a9) at (2,0) {9. Normativity};
\node[arg, fill=blue!20] (a10) at (8,0) {10. Vagueness};

% Connections
\draw[-] (epist) -- (a1);
\draw[-] (epist) -- (a2);
\draw[-] (epist) -- (a8);
\draw[-] (epist) -- (a9);
\draw[-] (epist) -- (a10);
\draw[-] (method) -- (a3);
\draw[-] (method) -- (a4);
\draw[-] (method) -- (a5);
\draw[-] (pract) -- (a6);
\draw[-] (pract) -- (a7);
\end{tikzpicture}
\caption{Classification of the Ten Canonical Arguments}
\label{fig:argument-classification}
\end{figure}

% ============================================================================
% ARGUMENT 1: IDENTIFICATION
% ============================================================================
\subsection{Argument 1: The Identification Problem}
\label{sec:arg-identification}

\begin{tcolorbox}[colback=red!5!white,colframe=red!75!black,title=\textbf{Core Claim}]
\textbf{``Complementarity is not uniquely identifiable from observational data.''}

Observed superadditivity can arise from:
\begin{itemize}[noitemsep]
    \item Latent common causes
    \item Selection processes
    \item Confounding variables
    \item Measurement artifacts
\end{itemize}
without genuine interaction.
\end{tcolorbox}

\subsubsection{Disciplinary Origins}

\begin{itemize}[noitemsep]
    \item \textbf{Economics:} Endogeneity in practice bundles (Athey \& Stern, 1998)
    \item \textbf{Epidemiology:} Confounding in interaction effects (VanderWeele, 2015)
    \item \textbf{Ecology:} Selection vs. complementarity (Loreau \& Hector, 2001)
    \item \textbf{Neuroscience:} FC $\neq$ effective connectivity (Friston, 2011)
\end{itemize}

\subsubsection{Argument Strength}

\textbf{Very High.} This is the most fundamental objection. The distinction between correlation and causation applies with special force to interactions, because:
\[
\text{Cov}(A \cdot B, Y) \neq \gamma_{AB} \cdot \text{Var}(A \cdot B)
\]

Any latent variable $Z$ that affects both $A$ and $B$ can create spurious interaction.

\subsubsection{Where the Argument Overreaches}

The identification problem is \textit{general}---it applies to all causal claims, not specifically to complementarity. The existence of identification challenges does not imply:
\begin{itemize}[noitemsep]
    \item Complementarity doesn't exist
    \item It cannot be estimated under appropriate designs
    \item Observational estimates are worthless
\end{itemize}

\subsubsection{$\gamma$-Framework Response}

\begin{enumerate}
    \item \textbf{Explicit decomposition:} The Möbius inversion makes the interaction term analytically separable
    \item \textbf{Design requirements:} Factorial experiments with randomization identify $\gamma$ causally
    \item \textbf{Sensitivity analysis:} Bounds under confounding (following Manski, 1990)
    \item \textbf{Instrumental strategies:} Natural experiments for quasi-experimental identification
\end{enumerate}

\subsubsection{Key References}

\textbf{Supporting the Objection:}
\begin{itemize}[noitemsep]
    \item Athey, S. \& Stern, S. (1998). Empirical framework for testing complementarities. \textit{NBER Working Paper}.
    \item VanderWeele, T. J. (2015). \textit{Explanation in Causal Inference}. Oxford.
    \item Imbens, G. W. \& Rubin, D. B. (2015). \textit{Causal Inference}. Cambridge.
\end{itemize}

\textbf{Addressing the Objection:}
\begin{itemize}[noitemsep]
    \item Milgrom, P. \& Roberts, J. (1995). Complementarities and fit. \textit{JAE}.
    \item Arora, A. \& Gambardella, A. (1990). Complementarity and external linkages. \textit{JIE}.
\end{itemize}

% ============================================================================
% ARGUMENT 2: MODEL-DEPENDENCE
% ============================================================================
\subsection{Argument 2: The Model-Dependence Problem}
\label{sec:arg-model-dependence}

\begin{tcolorbox}[colback=red!5!white,colframe=red!75!black,title=\textbf{Core Claim}]
\textbf{``Whether something is complementary depends on the reference model.''}

The same data can yield:
\begin{itemize}[noitemsep]
    \item $\gamma > 0$ under one reference model
    \item $\gamma < 0$ under another
    \item $\gamma = 0$ under a third
\end{itemize}
\end{tcolorbox}

\subsubsection{Canonical Examples}

\textbf{Pharmacology:} The classic case.

\begin{table}[h]
\centering
\caption{Drug Synergy Under Different Reference Models}
\label{tab:drug-synergy-models}
\renewcommand{\arraystretch}{1.3}
\begin{tabular}{@{}llcc@{}}
\toprule
\textbf{Model} & \textbf{Reference} & \textbf{Formula} & \textbf{Same Data} \\
\midrule
Loewe Additivity & Loewe (1928) & Isobologram & Synergy \\
Bliss Independence & Bliss (1939) & $E_{AB} = E_A + E_B - E_A E_B$ & Additivity \\
HSA & Berenbaum (1989) & $E_{AB} > \max(E_A, E_B)$ & Antagonism \\
ZIP & Yadav et al. (2015) & Zero interaction potency & Synergy \\
\bottomrule
\end{tabular}
\end{table}

\textbf{Machine Learning:} SHAP baseline dependence.

\begin{itemize}[noitemsep]
    \item Marginal baseline: Different $\phi$ values
    \item Interventional baseline: Different $\phi$ values
    \item Conditional baseline: Different $\phi$ values
\end{itemize}

\subsubsection{Argument Strength}

\textbf{High.} This is empirically well-documented. Meyer et al. (2019) show that different synergy metrics agree in only 30-50\% of cases.

\subsubsection{Where the Argument Overreaches}

Model-dependence is an argument against \textit{implicit} or \textit{unstated} reference models, not against complementarity \textit{per se}. The solution is:
\begin{itemize}[noitemsep]
    \item Make the reference model explicit
    \item Report results under multiple references
    \item Develop theory-driven reference selection
\end{itemize}

\subsubsection{$\gamma$-Framework Response}

The EBF approach:
\begin{enumerate}
    \item \textbf{Explicit null model:} $f(\emptyset)$ is always specified
    \item \textbf{Additivity as default:} $\gamma = 0$ means strict additivity
    \item \textbf{Sensitivity reporting:} Results under alternative references
    \item \textbf{Domain conventions:} Recommended references by field
\end{enumerate}

\subsubsection{Key References}

\textbf{Supporting the Objection:}
\begin{itemize}[noitemsep]
    \item Meyer, C. T. et al. (2019). Quantifying drug combination synergy... pitfalls. \textit{Trends Pharm Sci}.
    \item Greco, W. R. et al. (1995). The search for synergy: a critical review. \textit{Pharm Rev}.
    \item Sundararajan, M. \& Najmi, A. (2020). The many Shapley values. \textit{ICML}.
\end{itemize}

\textbf{Addressing the Objection:}
\begin{itemize}[noitemsep]
    \item Chou, T.-C. (2006). Theoretical basis... of synergism. \textit{Pharm Rev}.
    \item Tang, J. et al. (2015). What is synergy? A unified framework. \textit{Front Pharm}.
\end{itemize}

% ============================================================================
% ARGUMENT 3: SCALE
% ============================================================================
\subsection{Argument 3: The Scale Problem}
\label{sec:arg-scale}

\begin{tcolorbox}[colback=red!5!white,colframe=red!75!black,title=\textbf{Core Claim}]
\textbf{``Complementarity is not scale-invariant.''}

\begin{itemize}[noitemsep]
    \item Micro level: Strong synergies
    \item Meso level: Moderate effects
    \item Macro level: Saturation, neutralization
\end{itemize}
\end{tcolorbox}

\subsubsection{Empirical Evidence}

\textbf{Ecology:} Plot-level complementarity often disappears at landscape scale.

\textbf{Neuroscience:} Local connectivity patterns don't predict global network properties.

\textbf{Organizations:} Team-level complementarity doesn't aggregate to firm-level.

\subsubsection{Argument Strength}

\textbf{High.} Scale-dependence is empirically robust across many domains.

\subsubsection{Where the Argument Overreaches}

Scale-dependence $\neq$ non-existence. Many physical laws are scale-dependent (e.g., elasticity, viscosity). The appropriate response is:
\begin{itemize}[noitemsep]
    \item Multi-scale models
    \item Cross-level interactions
    \item Explicit aggregation rules
\end{itemize}

\subsubsection{$\gamma$-Framework Response}

The hierarchical extension (Appendix MDL):
\[
\gamma^L_{ab} = U^L(\{a,b\}) - U^L(\{a\}) - U^L(\{b\}) + U^L(\emptyset)
\]

Level-specific complementarity with explicit cross-level functions $\Lambda: \gamma^{L_1} \to \gamma^{L_2}$.

\subsubsection{Key References}

\textbf{Supporting the Objection:}
\begin{itemize}[noitemsep]
    \item Levin, S. A. (1992). The problem of pattern and scale in ecology. \textit{Ecology}.
    \item Wardle, D. A. (2016). Do experiments exploring plant diversity-ecosystem functioning tell us anything? \textit{Oikos}.
\end{itemize}

% ============================================================================
% ARGUMENT 4: CONTEXT
% ============================================================================
\subsection{Argument 4: The Context Problem}
\label{sec:arg-context}

\begin{tcolorbox}[colback=red!5!white,colframe=red!75!black,title=\textbf{Core Claim}]
\textbf{``Complementarity is unstable across contexts.''}

The same combination shows:
\begin{itemize}[noitemsep]
    \item Synergy in context $\Psi_1$
    \item Additivity in context $\Psi_2$
    \item Antagonism in context $\Psi_3$
\end{itemize}
\end{tcolorbox}

\subsubsection{Examples}

\begin{itemize}[noitemsep]
    \item \textbf{Drug combinations:} Synergy in vitro, additivity in vivo
    \item \textbf{Ecological:} Complementarity in benign conditions, not under stress
    \item \textbf{Organizational:} Practice complementarity in some industries, not others
    \item \textbf{Policy:} Intervention bundles work in some countries, not others
\end{itemize}

\subsubsection{Argument Strength}

\textbf{High.} This is the reproducibility crisis applied to interactions.

\subsubsection{Where the Argument Overreaches}

Context-dependence is \textit{expected} for complex phenomena. Many established effects are context-dependent:
\begin{itemize}[noitemsep]
    \item Price elasticities vary by market
    \item Treatment effects are heterogeneous
    \item Physical constants depend on conditions
\end{itemize}

\subsubsection{$\gamma$-Framework Response}

Conditional complementarity:
\[
\gamma_{ab}|\Psi = f(\{a,b\}|\Psi) - f(\{a\}|\Psi) - f(\{b\}|\Psi) + f(\emptyset|\Psi)
\]

The Context Framework (Appendix CTW) explicitly models $\Psi$.

\subsubsection{Key References}

\textbf{Supporting the Objection:}
\begin{itemize}[noitemsep]
    \item Lehar, J. et al. (2009). Context-dependent drug synergy. \textit{Nat Biotech}.
    \item Isbell, F. et al. (2015). Biodiversity increases the resistance of ecosystem productivity to climate extremes. \textit{Nature}.
\end{itemize}

% ============================================================================
% ARGUMENT 5: DOMINANCE/MASS-RATIO
% ============================================================================
\subsection{Argument 5: The Dominance Problem}
\label{sec:arg-dominance}

\begin{tcolorbox}[colback=red!5!white,colframe=red!75!black,title=\textbf{Core Claim}]
\textbf{``Effects come from dominant components, not from interaction.''}

Apparent complementarity is often:
\begin{itemize}[noitemsep]
    \item Selection effects (best component dominates)
    \item Mass-ratio effects (abundant component drives outcome)
    \item ``One best practice'' (single factor sufficient)
\end{itemize}
\end{tcolorbox}

\subsubsection{The BEF Debate}

This argument is central to the Biodiversity-Ecosystem Function literature:

\begin{quote}
``High-diversity plots often perform well because they are more likely to contain a high-performing species, not because species interact synergistically.'' ---Huston (1997)
\end{quote}

\textbf{Selection Effect:}
\[
\text{Performance} = \text{max}(\text{Species}_i) \quad \text{not} \quad \sum_i \gamma_{ij}
\]

\textbf{Mass-Ratio Effect (Grime, 1998):}
\[
\text{Function} = \sum_i p_i \cdot \text{Trait}_i
\]

where $p_i$ is relative abundance.

\subsubsection{Argument Strength}

\textbf{Medium-High.} Empirically supported in many BEF experiments (Smith et al., 2020).

\subsubsection{Where the Argument Overreaches}

Selection and complementarity are not mutually exclusive:
\begin{itemize}[noitemsep]
    \item Both can operate simultaneously
    \item Loreau \& Hector (2001) provide decomposition
    \item Complementarity often emerges over time
\end{itemize}

\subsubsection{$\gamma$-Framework Response}

Weighted decomposition:
\[
\Delta Y = \underbrace{\sum_i w_i \cdot \Delta \bar{Y}_i}_{\text{Selection}} + \underbrace{\sum_{i<j} \gamma_{ij} \cdot \Delta \bar{Y}_{ij}}_{\text{Complementarity}}
\]

Following Loreau \& Hector (2001) but with explicit $\gamma$.

\subsubsection{Key References}

\textbf{Supporting the Objection:}
\begin{itemize}[noitemsep]
    \item Grime, J. P. (1998). Benefits of plant diversity... immediate, filter, and founder effects. \textit{J Ecol}.
    \item Smith, M. D. et al. (2020). Mass ratio effects underlie ecosystem responses. \textit{J Ecol}.
    \item Huston, M. A. (1997). Hidden treatments in ecological experiments. \textit{Oecologia}.
\end{itemize}

\textbf{Addressing the Objection:}
\begin{itemize}[noitemsep]
    \item Loreau, M. \& Hector, A. (2001). Partitioning selection and complementarity. \textit{Nature}.
    \item Fox, J. W. (2005). Interpreting the ``selection effect'' of biodiversity. \textit{Ecol Lett}.
\end{itemize}

% ============================================================================
% ARGUMENT 6: COMPLEXITY
% ============================================================================
\subsection{Argument 6: The Complexity Problem}
\label{sec:arg-complexity}

\begin{tcolorbox}[colback=red!5!white,colframe=red!75!black,title=\textbf{Core Claim}]
\textbf{``Complementarity explodes combinatorially and is practically intractable.''}

For $n$ elements:
\[
\text{Interaction parameters} = 2^n - n - 1
\]

For $n = 20$: over 1 million parameters.
\end{tcolorbox}

\subsubsection{The Curse of Dimensionality}

\begin{table}[h]
\centering
\caption{Combinatorial Explosion of Interaction Terms}
\label{tab:combinatorial-explosion}
\renewcommand{\arraystretch}{1.2}
\begin{tabular}{@{}rrrr@{}}
\toprule
$n$ & Main Effects & Pairwise $\gamma$ & All Interactions \\
\midrule
5 & 5 & 10 & 26 \\
10 & 10 & 45 & 1,013 \\
20 & 20 & 190 & 1,048,555 \\
50 & 50 & 1,225 & $>10^{15}$ \\
\bottomrule
\end{tabular}
\end{table}

\subsubsection{Argument Strength}

\textbf{Very High.} This is a genuine computational and statistical barrier.

\subsubsection{Where the Argument Overreaches}

Complexity is a \textit{practical} limitation, not an \textit{ontological} one. Solutions exist:
\begin{itemize}[noitemsep]
    \item Sparsity assumptions (most $\gamma_{ij} \approx 0$)
    \item Hierarchical constraints (higher-order $\gamma$ bounded)
    \item Structural zeros (some combinations impossible)
    \item Approximation algorithms (polynomial-time bounds)
\end{itemize}

\subsubsection{$\gamma$-Framework Response}

\textbf{Möbius Structure (Appendix FND):}

The Möbius inversion provides exact decomposition:
\[
\gamma_S = \sum_{T \subseteq S} (-1)^{|S|-|T|} f(T)
\]

\textbf{Regularization:}
\begin{itemize}[noitemsep]
    \item LASSO for sparse $\gamma$
    \item Group penalties for structured sparsity
    \item Bayesian priors encoding decay with order
\end{itemize}

\textbf{Approximation:}
\begin{itemize}[noitemsep]
    \item Shapley sampling (polynomial)
    \item Kernel SHAP (efficient approximation)
    \item Monte Carlo estimation (Appendix LLM1)
\end{itemize}

\subsubsection{Key References}

\textbf{Supporting the Objection:}
\begin{itemize}[noitemsep]
    \item Weinreich, D. M. et al. (2013). Should evolutionary geneticists worry about higher-order epistasis? \textit{Curr Opin Genet Dev}.
\end{itemize}

\textbf{Addressing the Objection:}
\begin{itemize}[noitemsep]
    \item Nemhauser, G. et al. (1978). Approximations for submodular maximization. \textit{Math Prog}.
    \item Krause, A. \& Golovin, D. (2014). Submodular function maximization. \textit{Survey}.
    \item Lundberg, S. M. \& Lee, S.-I. (2017). SHAP. \textit{NeurIPS}.
\end{itemize}

% ============================================================================
% ARGUMENT 7: MEASUREMENT NOISE
% ============================================================================
\subsection{Argument 7: The Measurement Noise Problem}
\label{sec:arg-noise}

\begin{tcolorbox}[colback=red!5!white,colframe=red!75!black,title=\textbf{Core Claim}]
\textbf{``Interactions are especially sensitive to measurement noise.''}

Because $\gamma$ involves multiple terms:
\[
\gamma_{ab} = f(\{a,b\}) - f(\{a\}) - f(\{b\}) + f(\emptyset)
\]

errors in any term propagate and compound.
\end{tcolorbox}

\subsubsection{The Reliability Crisis in Neuroscience}

Noble et al. (2017, 2019) document test-retest reliability of functional connectivity:
\begin{itemize}[noitemsep]
    \item ICC often below 0.5
    \item High variability across sessions
    \item Motion artifacts create spurious connectivity
\end{itemize}

\subsubsection{Drug Screening Reliability}

High-throughput screens show:
\begin{itemize}[noitemsep]
    \item Batch effects larger than biological signal
    \item Lab-to-lab variation substantial
    \item ``Synergy'' often not reproducible
\end{itemize}

\subsubsection{Argument Strength}

\textbf{Very High.} This is empirically well-documented and practically important.

\subsubsection{Where the Argument Overreaches}

Measurement problems are arguments for:
\begin{itemize}[noitemsep]
    \item Better measurement (not abandoning the concept)
    \item Appropriate uncertainty quantification
    \item Replication requirements
\end{itemize}

\subsubsection{$\gamma$-Framework Response}

\begin{enumerate}
    \item \textbf{Error propagation formulas:} $\text{Var}(\gamma) = \sum_S c_S^2 \text{Var}(f(S))$
    \item \textbf{Regularization:} Shrinkage toward zero for noisy estimates
    \item \textbf{Bootstrap confidence intervals:} Honest uncertainty
    \item \textbf{Replication requirements:} Minimum studies for meta-analytic $\gamma$
\end{enumerate}

\subsubsection{Key References}

\textbf{Supporting the Objection:}
\begin{itemize}[noitemsep]
    \item Noble, S. et al. (2017). Test-retest reliability of functional connectivity. \textit{Cereb Cortex}.
    \item Noble, S. et al. (2019). A decade of test-retest reliability. \textit{NeuroImage}.
    \item Haibe-Kains, B. et al. (2013). Inconsistency in large pharmacogenomic studies. \textit{Nature}.
\end{itemize}

% ============================================================================
% ARGUMENT 8: POST-HOC NARRATIVITY
% ============================================================================
\subsection{Argument 8: The Post-Hoc Narrativity Problem}
\label{sec:arg-posthoc}

\begin{tcolorbox}[colback=red!5!white,colframe=red!75!black,title=\textbf{Core Claim}]
\textbf{``Complementarity is often a retrospective explanation, not a prediction.''}

\begin{itemize}[noitemsep]
    \item Successful systems appear ``well-integrated''
    \item Failures are rarely analyzed for ``missing complementarity''
    \item Survivorship bias inflates apparent complementarity
\end{itemize}
\end{tcolorbox}

\subsubsection{The Management Literature Problem}

In organizational studies:
\begin{itemize}[noitemsep]
    \item Toyota appears complementary because Toyota succeeded
    \item Failed companies aren't studied for ``anti-complementarity''
    \item ``Fit'' is diagnosed from outcomes, not measured independently
\end{itemize}

\subsubsection{Argument Strength}

\textbf{Medium.} This is a valid methodological concern but not unique to complementarity.

\subsubsection{Where the Argument Overreaches}

Post-hoc explanation is a general problem in social science. The solution is prospective design, not concept abandonment.

\subsubsection{$\gamma$-Framework Response}

\begin{enumerate}
    \item \textbf{Ex-ante specification:} $\gamma$ predictions before outcome observation
    \item \textbf{Pre-registration:} Registered reports with $\gamma$ hypotheses
    \item \textbf{Out-of-sample prediction:} Test $\gamma$ on new data
    \item \textbf{Failed combinations:} Explicitly study $\gamma < 0$ cases
\end{enumerate}

\subsubsection{Key References}

\textbf{Supporting the Objection:}
\begin{itemize}[noitemsep]
    \item Denrell, J. (2003). Vicarious learning, undersampling of failure. \textit{Org Sci}.
    \item Levinthal, D. A. (1997). Adaptation on rugged landscapes. \textit{Mgmt Sci}.
\end{itemize}

% ============================================================================
% ARGUMENT 9: NORMATIVITY
% ============================================================================
\subsection{Argument 9: The Normativity Problem}
\label{sec:arg-normativity}

\begin{tcolorbox}[colback=red!5!white,colframe=red!75!black,title=\textbf{Core Claim}]
\textbf{``Complementarity falsely suggests that `more is better'.''}

Risk of:
\begin{itemize}[noitemsep]
    \item Policy overload (too many simultaneous interventions)
    \item Polypharmacy (too many drugs)
    \item Overengineering (unnecessary complexity)
\end{itemize}
\end{tcolorbox}

\subsubsection{Real-World Risks}

\begin{itemize}[noitemsep]
    \item \textbf{Medicine:} Polypharmacy harms exceed benefits in elderly
    \item \textbf{Policy:} ``Christmas tree'' bills with contradictory provisions
    \item \textbf{Organizations:} Practice overload reduces performance
\end{itemize}

\subsubsection{Argument Strength}

\textbf{Medium.} This is a valid concern about misapplication.

\subsubsection{Where the Argument Overreaches}

The argument confuses:
\begin{itemize}[noitemsep]
    \item $\gamma > 0$ (complementarity exists) with
    \item ``Always add more'' (prescriptive error)
\end{itemize}

The $\gamma$-framework explicitly includes $\gamma < 0$ (antagonism).

\subsubsection{$\gamma$-Framework Response}

\begin{enumerate}
    \item \textbf{Negative $\gamma$ explicit:} Antagonism is part of the theory
    \item \textbf{Optimal bundles:} Find combinations that maximize, not add
    \item \textbf{Diminishing returns:} Higher-order $\gamma$ often negative
    \item \textbf{Cost functions:} Include complexity costs in optimization
\end{enumerate}

% ============================================================================
% ARGUMENT 10: CONCEPTUAL VAGUENESS
% ============================================================================
\subsection{Argument 10: The Conceptual Vagueness Problem}
\label{sec:arg-vagueness}

\begin{tcolorbox}[colback=red!5!white,colframe=red!75!black,title=\textbf{Core Claim}]
\textbf{``Complementarity is a container concept without clear boundaries.''}

Conflated with:
\begin{itemize}[noitemsep]
    \item Interaction (statistical)
    \item Emergence (philosophical)
    \item Correlation (associational)
    \item Coexistence (ecological)
    \item Synergy (pharmacological)
    \item Superadditivity (mathematical)
\end{itemize}
\end{tcolorbox}

\subsubsection{The Babel Problem}

As documented in Appendix HTY (LIT-HISTORY), the same mathematical structure is called:
\begin{itemize}[noitemsep]
    \item Complementarity (economics)
    \item Epistasis (genetics)
    \item Synergy (pharmacology)
    \item Interaction (statistics)
    \item Functional connectivity (neuroscience)
\end{itemize}

\subsubsection{Argument Strength}

\textbf{Very High.} This is the deepest philosophical objection.

\subsubsection{Where the Argument Overreaches}

Conceptual vagueness is the \textit{starting point} for formalization, not a refutation. The response is: provide precise definitions.

\subsubsection{$\gamma$-Framework Response}

This is exactly what the EBF framework does:

\begin{definition}[Formal Complementarity]
\[
\gamma_{ab} = f(\{a,b\}) - f(\{a\}) - f(\{b\}) + f(\emptyset)
\]
where $f: 2^S \to \mathbb{R}$ is a set function on elements $S$.
\end{definition}

This definition:
\begin{itemize}[noitemsep]
    \item Unifies all disciplinary variants
    \item Provides clear measurement procedures
    \item Distinguishes complementarity from related concepts
    \item Enables cross-domain comparison
\end{itemize}

\subsubsection{Key References}

\textbf{Supporting the Objection:}
\begin{itemize}[noitemsep]
    \item Grabisch, M. (2016). \textit{Set Functions, Games and Capacities}. Springer. [definitional chapter]
\end{itemize}

\textbf{Addressing the Objection:}
\begin{itemize}[noitemsep]
    \item Topkis, D. M. (1998). \textit{Supermodularity and Complementarity}. Princeton.
    \item Milgrom, P. \& Shannon, C. (1994). Monotone comparative statics. \textit{Econometrica}.
\end{itemize}

% ============================================================================
% PART II: LITERATURE DOCUMENTATION
% ============================================================================
\section{Literature Documentation: Pro and Contra}
\label{sec:literature}

\subsection{List A: 50 Papers Supporting Complementarity}

\subsubsection{Economics and Organizational Design (10 papers)}

\begin{enumerate}
    \item Milgrom, P. \& Roberts, J. (1990). Rationalizability, learning, and equilibrium in games with strategic complementarities. \textit{Econometrica}, 58(6), 1255--1277.

    \item Topkis, D. M. (1998). \textit{Supermodularity and Complementarity}. Princeton University Press.

    \item Athey, S. (2002). Monotone comparative statics under uncertainty. \textit{QJE}, 117(1), 187--223.

    \item Vives, X. (1990). Nash equilibrium with strategic complementarities. \textit{J Math Econ}, 19(3), 305--321.

    \item Bulow, J., Geanakoplos, J., \& Klemperer, P. (1985). Multimarket oligopoly: Strategic substitutes and complements. \textit{JPE}, 93(3), 488--511.

    \item Holmström, B. \& Milgrom, P. (1991). Multitask principal-agent analyses. \textit{JLEO}, 7, 24--52.

    \item Milgrom, P. \& Roberts, J. (1995). Complementarities and fit: Strategy, structure, and organizational change. \textit{JAE}, 19(2--3), 179--208.

    \item Brynjolfsson, E. \& Milgrom, P. (2013). Complementarity in organizations. In \textit{Handbook of Organizational Economics}. Princeton.

    \item Ichniowski, C., Shaw, K., \& Prennushi, G. (1997). The effects of human resource management practices on productivity. \textit{AER}, 87(3), 291--313.

    \item MacDuffie, J. P. (1995). Human resource bundles and manufacturing performance. \textit{ILR Review}, 48(2), 197--221.
\end{enumerate}

\subsubsection{Pharmacology and Medicine (10 papers)}

\begin{enumerate}
    \setcounter{enumi}{10}
    \item Chou, T.-C. \& Talalay, P. (1984). Quantitative analysis of dose-effect relationships. \textit{Adv Enzyme Regul}, 22, 27--55.

    \item Chou, T.-C. (2006). Theoretical basis, experimental design, and computerized simulation of synergism. \textit{Pharm Rev}, 58(3), 621--681.

    \item Berenbaum, M. C. (1989). What is synergy? \textit{Pharm Rev}, 41(2), 93--141.

    \item Loewe, S. (1953). The problem of synergism and antagonism of combined drugs. \textit{Arzneimittel-Forschung}, 3, 285--290.

    \item Bliss, C. I. (1939). The toxicity of poisons applied jointly. \textit{Ann Appl Biol}, 26, 585--615.

    \item Greco, W. R., Bravo, G., \& Parsons, J. C. (1995). The search for synergy: A critical review. \textit{Pharm Rev}, 47(2), 331--385.

    \item Lehar, J. et al. (2009). Synergistic drug combinations improve therapeutic selectivity. \textit{Nat Biotech}, 27(7), 659--666.

    \item Borisy, A. A. et al. (2003). Systematic discovery of multicomponent therapeutics. \textit{PNAS}, 100(13), 7977--7982.

    \item Fitzgerald, J. B. et al. (2006). Systems biology and combination therapy. \textit{Nat Chem Biol}, 2(9), 458--466.

    \item Ianevski, A. et al. (2022). SynergyFinder 3.0: Consensus interpretation of drug combinations. \textit{NAR}, 50(W1), W739--W743.
\end{enumerate}

\subsubsection{Ecology and Biodiversity (10 papers)}

\begin{enumerate}
    \setcounter{enumi}{20}
    \item Loreau, M. \& Hector, A. (2001). Partitioning selection and complementarity in biodiversity experiments. \textit{Nature}, 412, 72--76.

    \item Tilman, D. et al. (1997). The influence of functional diversity and composition on ecosystem processes. \textit{Science}, 277, 1300--1302.

    \item Hector, A. et al. (1999). Plant diversity and productivity experiments in European grasslands. \textit{Science}, 286, 1123--1127.

    \item Cardinale, B. J. et al. (2006). Effects of biodiversity on the functioning of trophic groups and ecosystems. \textit{Nature}, 443, 989--992.

    \item Cardinale, B. J. et al. (2012). Biodiversity loss and its impact on humanity. \textit{Nature}, 486, 59--67.

    \item Isbell, F. et al. (2011). High plant diversity is needed to maintain ecosystem services. \textit{Nature}, 477, 199--202.

    \item Reich, P. B. et al. (2012). Impacts of biodiversity loss escalate through time. \textit{Science}, 336, 589--592.

    \item Loreau, M. (2010). \textit{From Populations to Ecosystems: Theoretical Foundations for a New Ecological Synthesis}. Princeton.

    \item Tilman, D., Reich, P. B., \& Knops, J. (2006). Biodiversity and ecosystem stability in a decade-long grassland experiment. \textit{Nature}, 441, 629--632.

    \item Hooper, D. U. et al. (2005). Effects of biodiversity on ecosystem functioning: A consensus of current knowledge. \textit{Ecol Monogr}, 75(1), 3--35.
\end{enumerate}

\subsubsection{Neuroscience (10 papers)}

\begin{enumerate}
    \setcounter{enumi}{30}
    \item Friston, K. J. et al. (2003). Dynamic causal modelling. \textit{NeuroImage}, 19(4), 1273--1302.

    \item Friston, K. J. (2011). Functional and effective connectivity: A review. \textit{Brain Connectivity}, 1(1), 13--36.

    \item Tononi, G., Sporns, O., \& Edelman, G. M. (1994). A measure for brain complexity: Relating functional segregation and integration. \textit{PNAS}, 91(11), 5033--5037.

    \item Sporns, O. (2011). The human connectome: A complex network. \textit{Ann NY Acad Sci}, 1224(1), 109--125.

    \item Sporns, O. (2013). Network attributes for segregation and integration in the human brain. \textit{Curr Opin Neurobiol}, 23(2), 162--171.

    \item Bassett, D. S. \& Sporns, O. (2017). Network neuroscience. \textit{Nat Neurosci}, 20(3), 353--364.

    \item Deco, G. et al. (2015). Rethinking segregation and integration: Contributions of whole-brain modelling. \textit{Nat Rev Neurosci}, 16(7), 430--439.

    \item Honey, C. J. et al. (2007). Network structure of cerebral cortex shapes functional connectivity. \textit{PNAS}, 104(24), 10240--10245.

    \item Breakspear, M. (2017). Dynamic models of large-scale brain activity. \textit{Nat Neurosci}, 20(3), 340--352.

    \item Park, H.-J. \& Friston, K. (2013). Structural and functional brain networks: From connections to cognition. \textit{Science}, 342(6158), 1238411.
\end{enumerate}

\subsubsection{Machine Learning and Computer Science (10 papers)}

\begin{enumerate}
    \setcounter{enumi}{40}
    \item Shapley, L. S. (1953). A value for n-person games. In \textit{Contributions to the Theory of Games II}. Princeton.

    \item Lundberg, S. M. \& Lee, S.-I. (2017). A unified approach to interpreting model predictions. \textit{NeurIPS}, 4765--4774.

    \item Grabisch, M. \& Roubens, M. (1999). An axiomatic approach to the concept of interaction. \textit{Fuzzy Sets Syst}, 104(2), 181--196.

    \item Sundararajan, M., Taly, A., \& Yan, Q. (2017). Axiomatic attribution for deep networks. \textit{ICML}, 3319--3328.

    \item Tsai, C.-F. et al. (2023). Faith-Shap: The faithful Shapley interaction index. \textit{JMLR}, 24, 1--42.

    \item Agarwal, A. et al. (2019). Shapley Taylor interaction index. \textit{arXiv:1909.03200}.

    \item Nemhauser, G. L., Wolsey, L. A., \& Fisher, M. L. (1978). An analysis of approximations for maximizing submodular set functions. \textit{Math Prog}, 14(1), 265--294.

    \item Krause, A. \& Golovin, D. (2014). Submodular function maximization. In \textit{Tractability}. Cambridge.

    \item Friedman, J. H. \& Popescu, B. E. (2008). Predictive learning via rule ensembles. \textit{Ann Appl Stat}, 2(3), 916--954.

    \item Lou, Y. et al. (2013). Accurate intelligible models with pairwise interactions. \textit{KDD}, 623--631.
\end{enumerate}

\subsection{List B: 50 Papers Contrasting/Challenging Complementarity}

\subsubsection{Pharmacology: Model-Dependence and Pitfalls (10 papers)}

\begin{enumerate}
    \item Meyer, C. T. et al. (2019). Quantifying drug combination synergy... agreements, disagreements and pitfalls. \textit{Trends Pharm Sci}, 40(5), 290--303.

    \item Di Veroli, G. Y. et al. (2016). Combenefit: An interactive platform for analysis of drug combinations. \textit{Bioinformatics}, 32(18), 2866--2868. [Comparison showing divergent results]

    \item Foucquier, J. \& Guedj, M. (2015). Analysis of drug combinations: Current methodological landscape. \textit{Pharmacol Res Perspect}, 3(3), e00149.

    \item Tang, J. et al. (2015). What is synergy? A unified framework. \textit{Front Pharmacol}, 6, 181. [Documents model disagreements]

    \item Geary, N. (2013). Understanding synergy. \textit{Am J Physiol Endocrinol Metab}, 304(3), E237--E253. [Critique of overuse]

    \item Baeder, D. Y. et al. (2016). Antimicrobial combinations: Bliss independence and Loewe additivity. \textit{Antimicrob Agents Chemother}, 60(7), 4414--4423.

    \item Goldoni, M. \& Johansson, C. (2007). A mathematical approach to study combined effects. \textit{Arch Toxicol}, 81(2), 83--91.

    \item Yeh, P. J. et al. (2009). Drug interactions and the evolution of antibiotic resistance. \textit{Nat Rev Microbiol}, 7(6), 460--466. [Context-dependence]

    \item Russ, D. \& Bhowmick, S. (2013). Synergy assessment... lack of standardization. \textit{Pharmacol Res}, 72, 34--48.

    \item Lederer, S. et al. (2018). The need for a standardized synergy score. \textit{Clin Pharmacol Ther}, 104(5), 844--847.
\end{enumerate}

\subsubsection{Ecology: Selection vs. Complementarity Debate (10 papers)}

\begin{enumerate}
    \setcounter{enumi}{10}
    \item Grime, J. P. (1998). Benefits of plant diversity to ecosystems: Immediate, filter, and founder effects. \textit{J Ecol}, 86(6), 902--910.

    \item Smith, M. D. et al. (2020). Mass ratio effects underlie ecosystem responses to environmental change. \textit{J Ecol}, 108(3), 855--864.

    \item Huston, M. A. (1997). Hidden treatments in ecological experiments. \textit{Oecologia}, 110(4), 449--460.

    \item Wardle, D. A. (1999). Is ``sampling effect'' a problem for experiments investigating biodiversity-ecosystem function? \textit{Oikos}, 87(2), 403--407.

    \item Thompson, K. et al. (2005). Do plant traits influence mechanisms of ecosystem function? \textit{Oikos}, 111(3), 624--628.

    \item Naeem, S. (2002). Ecosystem consequences of biodiversity loss: The evolution of a paradigm. \textit{Ecology}, 83(6), 1537--1552. [Critical review]

    \item Leps, J. et al. (2006). How reliable are our vegetation analyses? \textit{J Veg Sci}, 17(3), 387--398.

    \item Schwartz, M. W. et al. (2000). Linking biodiversity to ecosystem function: Implications for conservation. \textit{Oecologia}, 122(3), 297--305.

    \item Srivastava, D. S. \& Vellend, M. (2005). Biodiversity-ecosystem function research: Is it relevant to conservation? \textit{Ann Rev Ecol Evol Syst}, 36, 267--294.

    \item Balvanera, P. et al. (2006). Quantifying the evidence for biodiversity effects on ecosystem functioning. \textit{Ecol Lett}, 9(10), 1146--1156. [Meta-analysis showing heterogeneity]
\end{enumerate}

\subsubsection{Neuroscience: Reliability and Validity Problems (10 papers)}

\begin{enumerate}
    \setcounter{enumi}{20}
    \item Noble, S. et al. (2017). Influences on the test-retest reliability of functional connectivity MRI. \textit{Cereb Cortex}, 27(11), 5415--5429.

    \item Noble, S. et al. (2019). A decade of test-retest reliability of functional connectivity: A systematic review. \textit{NeuroImage}, 203, 116157.

    \item Birn, R. M. et al. (2013). The effect of scan length on the reliability of resting-state fMRI connectivity estimates. \textit{NeuroImage}, 83, 550--558.

    \item Laumann, T. O. et al. (2015). Functional system and areal organization of a highly sampled individual human brain. \textit{Neuron}, 87(3), 657--670. [Individual variability]

    \item Power, J. D. et al. (2012). Spurious but systematic correlations in functional connectivity MRI networks arise from subject motion. \textit{NeuroImage}, 59(3), 2142--2154.

    \item Van Dijk, K. R. A. et al. (2012). The influence of head motion on intrinsic functional connectivity MRI. \textit{NeuroImage}, 59(1), 431--438.

    \item Poldrack, R. A. et al. (2017). Scanning the horizon: Towards transparent and reproducible neuroimaging research. \textit{Nat Rev Neurosci}, 18(2), 115--126.

    \item Button, K. S. et al. (2013). Power failure: Why small sample size undermines the reliability of neuroscience. \textit{Nat Rev Neurosci}, 14(5), 365--376.

    \item Eklund, A. et al. (2016). Cluster failure: Why fMRI inferences... are inflated. \textit{PNAS}, 113(28), 7900--7905.

    \item Marek, S. et al. (2022). Reproducible brain-wide association studies require thousands of individuals. \textit{Nature}, 603, 654--660.
\end{enumerate}

\subsubsection{ML/XAI: Causal Interpretation Problems (10 papers)}

\begin{enumerate}
    \setcounter{enumi}{30}
    \item Janzing, D. et al. (2020). Feature relevance quantification in explainable AI: A causal problem. \textit{AISTATS}, 2907--2916.

    \item Sundararajan, M. \& Najmi, A. (2020). The many Shapley values for model explanation. \textit{ICML}, 9269--9278.

    \item Kumar, I. E. et al. (2020). Problems with Shapley-value-based explanations as feature importance measures. \textit{ICML}, 5491--5500.

    \item Slack, D. et al. (2020). Fooling LIME and SHAP: Adversarial attacks on post hoc explanation methods. \textit{AIES}, 180--186.

    \item Rudin, C. (2019). Stop explaining black box machine learning models for high stakes decisions. \textit{Nat Mach Intell}, 1(5), 206--215.

    \item Hooker, S. et al. (2019). A benchmark for interpretability methods in deep neural networks. \textit{NeurIPS}, 9737--9748.

    \item Aas, K. et al. (2021). Explaining individual predictions when features are dependent. \textit{Artif Intell}, 298, 103502.

    \item Chen, H. et al. (2020). True to the model or true to the data? \textit{arXiv:2006.16234}.

    \item Merrick, L. \& Taly, A. (2020). The explanation game: Explaining machine learning models using Shapley values. \textit{arXiv:2003.05316}.

    \item Frye, C. et al. (2021). Shapley explainability on the data manifold. \textit{ICLR}.
\end{enumerate}

\subsubsection{Economics/Social Science: Identification and Endogeneity (10 papers)}

\begin{enumerate}
    \setcounter{enumi}{40}
    \item Athey, S. \& Stern, S. (1998). An empirical framework for testing theories about complementarity in organizational design. \textit{NBER Working Paper} 6600.

    \item Arora, A. (1996). Testing for complementarities in reduced-form regressions: A note. \textit{Econ Lett}, 50(1), 51--55.

    \item Cassiman, B. \& Veugelers, R. (2006). In search of complementarity in innovation strategy. \textit{Mgmt Sci}, 52(1), 68--82. [Identification challenges]

    \item Mohnen, P. \& Röller, L.-H. (2005). Complementarities in innovation policy. \textit{Eur Econ Rev}, 49(6), 1431--1450.

    \item Miravete, E. J. \& Pernías, J. C. (2006). Innovation complementarity and scale of production. \textit{J Ind Econ}, 54(1), 1--29.

    \item Bloom, N. et al. (2012). Americans do IT better. \textit{AER}, 102(1), 167--201. [Complementarity hard to identify]

    \item Leiponen, A. (2005). Organization of knowledge and innovation: The case of Finnish business services. \textit{Ind Innov}, 12(2), 185--203.

    \item Bocquet, R. et al. (2007). Complementarities in organizational design. \textit{Ind Corp Change}, 16(4), 521--549.

    \item Ennen, E. \& Richter, A. (2010). The whole is more than the sum of its parts—or is it? A review of the empirical literature on complementarities. \textit{J Mgmt}, 36(1), 207--233.

    \item Kok, R. A. W. \& Creemers, M. (2008). Alliance governance and product innovation project decision making. \textit{Eur J Innovation Mgmt}, 11(4), 472--487.
\end{enumerate}

% ============================================================================
% PART III: META-ANALYSIS
% ============================================================================
\section{Meta-Analysis: The Structure of Criticism}
\label{sec:meta-analysis}

\subsection{Classification of Objections}

\begin{table}[h]
\centering
\caption{Classification of the Ten Arguments}
\label{tab:argument-classification}
\renewcommand{\arraystretch}{1.3}
\begin{tabular}{@{}clccc@{}}
\toprule
\textbf{\#} & \textbf{Argument} & \textbf{Type} & \textbf{Ontological?} & \textbf{Solvable?} \\
\midrule
1 & Identification & Epistemological & No & Yes \\
2 & Model-Dependence & Methodological & No & Yes \\
3 & Scale & Methodological & No & Yes \\
4 & Context & Methodological & No & Yes \\
5 & Dominance & Empirical & No & Yes \\
6 & Complexity & Practical & No & Yes \\
7 & Noise & Practical & No & Yes \\
8 & Post-hoc & Epistemological & No & Yes \\
9 & Normativity & Conceptual & No & Yes \\
10 & Vagueness & Conceptual & No & Yes \\
\bottomrule
\end{tabular}
\end{table}

\subsection{The Critical Observation}

\begin{tcolorbox}[colback=green!5!white,colframe=green!75!black,title=\textbf{Central Finding}]
\textbf{None of the ten arguments is ontological.}

All ten are:
\begin{itemize}[noitemsep]
    \item Methodological (how to measure)
    \item Epistemological (how to know)
    \item Practical (how to compute)
    \item Conceptual (how to define)
\end{itemize}

None claims that complementarity \textit{does not exist}.
\end{tcolorbox}

This is the key insight: the objections explain \textit{why complementarity was under-formalized}, not \textit{why it is false}.

\subsection{How the $\gamma$-Framework Addresses Each Argument}

\begin{table}[h]
\centering
\caption{$\gamma$-Framework Responses to Canonical Arguments}
\label{tab:gamma-responses}
\renewcommand{\arraystretch}{1.4}
\small
\begin{tabular}{@{}clp{6cm}@{}}
\toprule
\textbf{\#} & \textbf{Argument} & \textbf{$\gamma$-Framework Response} \\
\midrule
1 & Identification & Explicit Möbius decomposition; factorial designs; sensitivity analysis \\
2 & Model-Dependence & Reference model always explicit; additivity as default; sensitivity to alternatives \\
3 & Scale & Hierarchical $\gamma^L$ at multiple levels; cross-level functions $\Lambda$ \\
4 & Context & Conditional $\gamma|\Psi$; Context Framework (Appendix CTW) \\
5 & Dominance & Weighted decomposition separating selection and complementarity \\
6 & Complexity & Möbius structure; SHAP sampling; regularization for sparsity \\
7 & Noise & Error propagation formulas; bootstrap CI; regularized estimation \\
8 & Post-hoc & Ex-ante specification; pre-registration; out-of-sample testing \\
9 & Normativity & $\gamma < 0$ explicit; optimization includes cost; diminishing returns \\
10 & Vagueness & Formal definition: $\gamma_{ab} = f(\{a,b\}) - f(\{a\}) - f(\{b\}) + f(\emptyset)$ \\
\bottomrule
\end{tabular}
\end{table}

% ============================================================================
% SUMMARY
% ============================================================================

%%%%%%%%%%%%%%%%%%%%%%%%%%%%%%%%%%%%%%%%%%%%%%%%%%%%%%%%%%%%%%%%%%%%%%%%%%%%%%%
\section{Key Results}
\label{sec:crt-results}
%%%%%%%%%%%%%%%%%%%%%%%%%%%%%%%%%%%%%%%%%%%%%%%%%%%%%%%%%%%%%%%%%%%%%%%%%%%%%%%

The analysis yields the following key results:

\begin{enumerate}
    \item \textbf{Result 1:} [Primary finding with quantitative support]
    \item \textbf{Result 2:} [Secondary finding with implications]
    \item \textbf{Result 3:} [Practical application or boundary condition]
\end{enumerate}


\section{Summary}
\label{sec:summary}

\begin{tcolorbox}[colback=blue!5!white,colframe=blue!75!black,title=\textbf{Concluding Statement}]
\begin{center}
\large
\textbf{The central objections to complementarity are not substantive but methodological---and therefore solvable.}
\end{center}

\vspace{0.5em}
The ten canonical arguments against complementarity explain:
\begin{enumerate}
    \item Why complementarity remained under-developed for 2,500 years (Appendix HTY)
    \item What a formal theory must address to be scientifically rigorous
    \item Why the current moment enables what previous eras could not
\end{enumerate}

The EBF $\gamma$-framework is the systematic response to all ten objections.
\end{tcolorbox}

\subsection{Cross-Reference: The Eleventh Argument}

This appendix documents ten \textit{technical} objections to complementarity. For a
deeper \textit{meta-level} critique---not about measurement or identification, but
about whether integrative frameworks can be scientific at all---see \textbf{Appendix FRM
(LIT-FEHR-METHOD): Integration vs. Falsifiability}. That appendix addresses the
question: ``If a framework can integrate all models, does it explain anything?''
The answer: $\gamma$ is scientifically legitimate when it commits to $\gamma = 0$
null hypotheses, making complementarity an \textit{exclusion principle} rather than
an integration principle.

% ============================================================================
% GLOSSARY
% ============================================================================
\section{Glossary}
\label{sec:glossary}

Key terms (see Appendix G for comprehensive Master Glossary):

\begin{description}[style=nextline]
    \item[Canonical Argument] A recurrent, structured objection that appears across multiple disciplines
    \item[Identification Problem] Difficulty distinguishing true interaction from confounded association
    \item[Model-Dependence] Sensitivity of complementarity estimates to reference model choice
    \item[Selection Effect] Apparent complementarity due to sampling high-performing components
    \item[Mass-Ratio Effect] Outcome driven by abundant components, not by interaction
    \item[Möbius Inversion] Mathematical decomposition separating interaction terms exactly
    \item[SHAP] SHapley Additive exPlanations; ML interpretability method
\end{description}

% ============================================================================
% OPEN ISSUES
% ============================================================================
\section{Open Issues and Research Agenda}
\label{sec:crt-open-issues}

\begin{enumerate}
    \item \textbf{Empirical Validation:} Further testing of canonical argument resolutions across domains
    \item \textbf{New Critiques:} Monitoring for emerging objections in AI/ML literature
    \item \textbf{Synthesis:} Developing unified response framework across disciplines
\end{enumerate}

% ============================================================================
% REFERENCES SUMMARY
% ============================================================================
\section{Reference Summary}

\begin{table}[h]
\centering
\caption{Reference Distribution}
\label{tab:reference-summary}
\renewcommand{\arraystretch}{1.2}
\begin{tabular}{@{}lrr@{}}
\toprule
\textbf{Category} & \textbf{Pro} & \textbf{Contra} \\
\midrule
Economics/Organizations & 10 & 10 \\
Pharmacology/Medicine & 10 & 10 \\
Ecology/Biodiversity & 10 & 10 \\
Neuroscience & 10 & 10 \\
ML/Computer Science & 10 & 10 \\
\midrule
\textbf{Total} & \textbf{50} & \textbf{50} \\
\bottomrule
\end{tabular}
\end{table}

\nocite{bcm_master}

\end{chapter}


%%%%%%%%%%%%%%%%%%%%%%%%%%%%%%%%%%%%%%%%%%%%%%%%%%%%%%%%%%%%%%%%%%%%%%%%%%%%%%%
\section{References}
\label{sec:crt-references}
%%%%%%%%%%%%%%%%%%%%%%%%%%%%%%%%%%%%%%%%%%%%%%%%%%%%%%%%%%%%%%%%%%%%%%%%%%%%%%%

See master bibliography (\texttt{bcm\_master.bib}) for complete citations.

\nocite{bcm_master}

